\documentclass{article}

\usepackage[utf8]{inputenc}
\usepackage[spanish]{babel}

\title{Análisis de la Calidad del Aire en el Área Metropolitana de Monterrey}
\author{María Clarissa Contreras Acosta}
\date{\today}

\begin{document}

\maketitle

\section{Introducción}

La calidad del aire es un tema de interés para el análisis
ambiental de las zonas urbanas. En este proyecto se plantea
un análisis de las concentraciones de contaminantes atmosféricos
en el Área Metropolitana de Monterrey.

\section{Objetivo}

El objetivo es analizar y visualizar las concentraciones de
PM2.5 y PM10 mediante herramientas computacionales.

\section{Metodología}

Para el desarrollo del proyecto se utiliza Python para realizar
cálculos estadísticos y generar visualizaciones. Los datos
utilizados en esta etapa son simulados con fines de desarrollo
del proyecto.

\section{Resultados preliminares}

Se calcularon los valores promedio, mínimo y máximo de PM2.5
y PM10. También se generó una gráfica para visualizar las
concentraciones de ambos contaminantes.

\section{Conclusiones}

El proyecto establece una base para realizar posteriormente
un análisis más detallado de la calidad del aire mediante
datos reales.

\end{document}